% SPDX-FileCopyrightText: 2026 David Rae Wilmot
% SPDX-License-Identifier: AGPL-3.0-or-later
%
% Shadow-Loom: an experimental project overview paper, ACL style.
% Build with the official acl-style-files (https://github.com/acl-org/acl-style-files):
%   pdflatex shadow_loom
%   bibtex   shadow_loom
%   pdflatex shadow_loom
%   pdflatex shadow_loom

\documentclass[11pt]{article}

% ACL 2023+ style. Place acl.sty / acl_natbib.bst alongside this file.
\usepackage[final]{acl}

\usepackage{times}
\usepackage{latexsym}
\usepackage[T1]{fontenc}
\usepackage[utf8]{inputenc}
\usepackage{microtype}
\usepackage{amsmath}
\usepackage{amssymb}
\usepackage{booktabs}
\usepackage{graphicx}
\usepackage{url}
\usepackage{hyperref}

\title{Shadow-Loom: An Experimental Framework for Combining
       Graphical Causal Reasoning with a Game-Engine-Style
       Graphical World Model of Narrative}

\author{David R.\ Wilmot \\
        david.wilmot@gmail.com \\
        \texttt{\href{https://github.com/dwlmt/shadow-loom}{github.com/davidrwilmot/shadow-loom}} \\
        Released under AGPL-3.0-or-later + Non-Commercial Clause}

\begin{document}
\maketitle

\begin{abstract}
Stories interest us not because they are sequences of likely
next words, but because they have causes, secrets, and
consequences. \textbf{Shadow-Loom} is an experimental open-source
framework that takes a narrative and turns it into a typed,
versioned \emph{graphical world model}. Two engines act on that
model, in the spirit of a game engine that pairs a physics
simulator with a higher-level director. The first engine is a
\emph{causal physics} grounded in Pearl's ladder of causation
\citep{pearl2009causality,pearl2018book,bareinboim2022pearl} and
the recently proposed counterfactual calculus over Ancestral
Multi-World Networks \citep{correa2025amwn}. The second is a
\emph{narrative physics} that scores the same graph against four
structural reader-states --- mystery, dramatic irony, suspense
\citep{wilmot2020suspense}, and surprise \citep{ittibaldi2009surprise}
--- in the tradition of Sternberg's curiosity / suspense /
surprise triad \citep{sternberg1992telling}, and plays the role of
the drama manager in classical interactive-narrative architectures
\citep{mateas2003facade,riedl2013interactive,ware2014glaive}.
Large language models are used only at the boundary: extraction,
rendering, and audit. Identification, intervention, and
counterfactual reasoning happen in typed code over the graph, much
as a game engine restricts a stochastic agent to physically
admissible actions \citep{ha2018world,park2023generative}. The
same loop --- ingest, simulate, generate, audit, re-ingest ---
runs over user prose and produces a versioned world model with
explicit, inspectable mathematical structure.

Two motivations drive the design. First, counterfactual
reasoning sits exactly on the rung of Pearl's hierarchy where
current LLMs are weakest. They are competent at observational
and even interventional causal questions in text
\citep{kiciman2024causal}, but they fail to act on the implicit
beliefs and counterfactual states they themselves infer
\citep{gu2026simpletom,kim2023fantom,xu2025echoes}, and the
foundations of rung-3 reasoning need dedicated graphical machinery
\citep{correa2025amwn,bareinboim2022pearl,halpern2016actual}.
Second, coding agents reliably outperform open-domain agents at
fixed model capacity because their outputs can be compiled,
executed, and tested against a typed substrate. Shadow-Loom
gives narrative reasoning the same affordance: candidate
continuations are propagated through the graph, checked against
counterfactual-calculus-style independence and exclusion rules
\citep{correa2025amwn}, and re-ingested for round-trip verification
\citep{madaan2023selfrefine,bai2022constitutional,gu2024judge}.

The system is presented as a research artefact, not as a
benchmarked NLP model. Formal evaluation is left to future work.
The more modest claim made here is that wiring a
causal-graphical substrate underneath a constrained-LLM renderer
is a tractable design point, with plausible relevance to
controllable generation, narrative theory of mind, computational
social science, and the digital humanities. Code, fixtures, and
pipeline are released open source so that other researchers --- in
NLP, in computational social science, and in the digital
humanities --- can adapt the substrate to their own questions.
Definitions and equations for each pipeline stage are given in
the appendix.
\end{abstract}


\section{Introduction}
\label{sec:intro}

Most current LLM-driven story systems treat a narrative as a
sequence-prediction problem and absorb the causal, temporal, and
epistemic structure of the story into the weights of a single
model. Recent work shows that such models are competent at
\emph{reading} causal relations from text
\citep{kiciman2024causal} but markedly weaker at \emph{acting} on
the implicit beliefs of characters they have just inferred
\citep{gu2026simpletom,kim2023fantom,cross2024hyp}; in
unconditioned generation they tend to collapse onto a small set
of shared plot tropes \citep{xu2025echoes,tian2024human}. The
complementary direction explored here keeps an explicit, typed
graph of the storyworld
\citep{schank1977scripts,mani2012computational,elson2012sig} and
performs identification, intervention, and counterfactual
reasoning over it symbolically
\citep{pearl2009causality,halpern2016actual,correa2025amwn},
using the LLM only as a constrained renderer
\citep{riedlyoung2010narrative,yao2019planwrite,
goldfarbtarrant2020aristotelian,yang2023doc}.

\paragraph{Motivation 1: counterfactuals are exactly where LLMs
struggle.} Pearl's hierarchy is provably non-collapsible. Rung-3
quantities are not in general identifiable from rung-2 data, let
alone from rung-1 observation \citep{bareinboim2022pearl}.
Empirically, LLMs reproduce surface causal patterns from their
training corpus \citep{kiciman2024causal} but degrade sharply on
multi-hop belief-conditioned action prediction
\citep{kim2023fantom,gu2026simpletom} and on story-internal
counterfactual consistency \citep{xu2025echoes}. The questions
that matter to a reader of a story --- ``what if Macbeth had
refused?'', ``what if Roy had told Mary the truth?'' --- live on
the rung where the modern foundations (AMWN + ctf-calculus,
\citealp{correa2025amwn}) are designed to operate, and where
end-to-end neural generators are not. A typed graph plus an
explicit do-operator and abduction step turns each such question
into a finite sequence of well-defined operations
\citep{pearl2009causality,halpern2005causes,halpern2016actual}
rather than into a free-form imagination task.

\paragraph{Motivation 2: coding agents work better because they
compile and test.} A familiar pattern in agent design is that
coding agents outperform open-domain agents at fixed model
capacity. Their outputs can be compiled, executed against a test
suite, and rejected when they fail. The lever is the
\emph{typed substrate underneath the natural-language reasoning},
not the model itself. Shadow-Loom imports that lever into
narrative. A continuation is rejected if (a) its causal
propagation (Eq.~\ref{eq:impact}) violates inertia or affordance
gating; (b) its counterfactual branch fails an AMWN-style
consistency, independence, or exclusion check
\citep{correa2025amwn}; or (c) its re-extraction does not match
the brief that licensed it (the audit pass, in the LLM-as-judge
tradition of \citealp{bai2022constitutional,
madaan2023selfrefine,gu2024judge}). The graph plays the role of
the compiler; the auditor plays the role of the test suite.

The framework combines four ideas which, to the best of my
knowledge, have not been wired together in a single open-source
pipeline:

\begin{enumerate}
\item A \textbf{Pydantic-typed world graph} (\texttt{WorldStateV1})
      with entities, events, beliefs, locations, channels, and three
      kinds of edge (causal, social, spatial), every node and edge
      tagged with both a chronological \emph{fabula} time and a
      narrative \emph{syuzhet} index \citep{shklovsky1965art,
      genette1980narrative}.
\item A \textbf{causal physics engine} that implements Pearl's three
      rungs --- observation, do-intervention, and abductive
      counterfactual --- with a propagation rule gated by an
      explicit ``Impact $>$ Inertia'' threshold and spatial
      affordance checks.
\item A \textbf{narrative physics layer} that scores the same graph
      against closed-form information-state functionals for mystery,
      dramatic irony, suspense, and surprise. Suspense follows the
      hope/threat formulation of \citet{wilmot2020suspense};
      surprise is binary KL divergence between a corpus-marginal
      prior and an evidence-updated posterior.
\item A \textbf{constrained-LLM renderer + recursive auditor} loop
      \citep{gu2024judge} that converts the simulation result into
      a typed creative brief, generates prose under that brief, and
      then re-extracts the prose to verify no ``miracle steps''
      were introduced.
\end{enumerate}

The system is named for the dual-graph metaphor: a \emph{loom} of
factual edges runs in parallel with one or more \emph{shadow}
counterfactual branches, in the sense of the AMWN construction of
\citet{correa2025amwn}. Section~\ref{sec:pipeline} sketches the
pipeline; Section~\ref{sec:novelty} states what is novel and what
is borrowed; Section~\ref{sec:rationale} gives the design
rationale; Section~\ref{sec:relevance} discusses relevance to NLP,
reasoning, and social science; Section~\ref{sec:limits} states the
limitations honestly. Equations and per-stage definitions are in
the appendix.


\section{Pipeline overview}
\label{sec:pipeline}

The pipeline is a five-phase loop over a versioned world model.
Every version is ancestor-linked, so a counterfactual fork is
literally a sibling row in a directed acyclic version tree, not a
mutation of canon. Figure~\ref{fig:pipeline} sketches the loop.

\begin{figure}[h]
\centering
\small
\begin{tabular}{c}
\texttt{prose}\\
$\downarrow$ \scriptsize ingestion (5-pass extraction)\\
\texttt{WorldStateV1 (typed graph)}\\
$\downarrow$ \scriptsize ego-graph slice + AMWN sandbox\\
\texttt{NetworkX MultiDiGraph}\\
$\downarrow$ \scriptsize causal physics (do, abduction, propagation)\\
\texttt{candidate sandbox states}\\
$\downarrow$ \scriptsize narrative physics (affective scoring)\\
\texttt{CreativeBrief (typed constraints)}\\
$\downarrow$ \scriptsize constrained LLM render\\
\texttt{prose}\\
$\downarrow$ \scriptsize recursive audit + re-extraction\\
\texttt{new VersionRow (factual or shadow)}\\
\end{tabular}
\caption{The Shadow-Loom loop. The LLM is invoked only at the top
(ingestion), the bottom (rendering), and the audit; everything in
between is symbolic computation over the graph. See the appendix
for definitions of each stage.}
\label{fig:pipeline}
\end{figure}

The five phases correspond to the five appendix subsections
(\ref{app:ingestion}--\ref{app:audit}):
\textbf{(1) Ingestion} runs five LLM agents to extract a typed
\texttt{GlobalRegister}, then three concurrent specialist agents
(physics, social, consequences) per chunk to extract events,
relationships, channels, and entity-state deltas, with a
programmatic validation and auto-repair loop.
\textbf{(2) Causal physics} forks an Ancestral Multi-World Network
sandbox, applies a do-operator for interventions, runs abduction for
counterfactuals, and propagates effects through a topological sort
of the causal sub-graph subject to inertia and affordance gating.
\textbf{(3) Narrative physics / affective calculus} scores each
candidate sandbox against four structural information-state
functionals and six trait-trajectory emotional functionals.
\textbf{(4) Generation} packages the highest-scoring candidate as a
\texttt{CreativeBrief} of typed \texttt{ConstraintBlock}s and hands
\emph{only} the brief to a creative LLM.
\textbf{(5) Audit and feedback} runs an LLM-as-judge over the prose,
detects ``miracle steps'' (state changes with no licensing edge in
the brief), and either accepts the version or re-runs generation
with the auditor's feedback appended, up to a configurable retry
budget.


\section{Theoretical lineage and computational counterparts}
\label{sec:lineage}

The system has been built theory-first: every named field in
\texttt{WorldStateV1} corresponds to a specific narratological,
cognitive, or causal-inference construct, and every field has a
prior computational instantiation that the design either
generalises or borrows from. The pairings the design depends on
are summarised below; an extended mapping lives in the project's
academic-foundations document.

\paragraph{Fabula vs syuzhet
\citep{shklovsky1965art,tomashevsky1965thematics,
genette1980narrative,bordwell1985narration,bal2009narratology}.}
Two integer indices on every \texttt{EventNode}: chronological
$f(e)$ and presentation $s(e)$. The four affective scorers reduce
to set operations between the two projections (App.~\ref{app:narrative}).

\paragraph{Possible-worlds narratology
\citep{ryan1991possible,ryan2001virtual,dolezel1998heterocosmica,
pavel1986fictional}, modal-logical roots
\citep{lewis1973causation}, and counterfactual graphical models
\citep{balke1994counterfactual,shpitser2008complete,
correa2021nested,correa2025amwn}.} Every node and edge carries
$\texttt{world\_id} \in \{\text{factual}, \text{shadow}\}$;
counterfactual queries fork a sibling \texttt{VersionRow}.
The AMWN tag and three-rule mental model (consistency,
independence, exclusion) come from \citet{correa2025amwn}; I
adopt the naming and architecture even though the typed narrative
graph used here is not a rigorous SCM. AMWN supersedes the Twin-Network
construction \citep{balke1994counterfactual} and the multi-network
machinery of \citet{shpitser2008complete}, and is the
counterfactual-completeness counterpart to do-calculus
\citep{pearl1995causal,pearl2009causality,bareinboim2022pearl}.

\paragraph{Greimas \citep{greimas1983structural}, Propp
\citep{propp1968morphology}, Bremond \citep{bremond1973logique},
and the symbolic story-grammar tradition
\citep{rumelhart1975schema,mandlerjohnson1977,kintsch1978macro,
schank1977scripts,lehnert1981plotunits,meehan1977talespin,
perez2001mexica,bringsjord2000brutus,chambers2008narrative}.}
\texttt{GlobalTrait} (\texttt{WORLD\_}) realises Greimas's ``Power''
actant; \texttt{RelationshipEdge.power\_dynamic} mirrors the
sender / receiver / helper / opponent topology;
\texttt{EventNode.event\_type} $\in \{\text{choice, outcome,
revelation, utterance}\}$ is a coarse generalisation of Propp's
functions,
closer in spirit to Bremond's triad. The plot-unit and script
traditions prefigure both the typed-event graph
\citep{schank1977scripts,lehnert1981plotunits} and the
goal-directed planner / generator pipeline
\citep{meehan1977talespin,perez2001mexica}; the unsupervised
narrative event chains of \citet{chambers2008narrative} are a direct
precursor to the typed \texttt{causal\_topology} maintained at
runtime.

\paragraph{Cognitive narratology
\citep{herman2002story,fludernik1996natural,zwaan1998situation,
gerrig1993narrative} and constructionist comprehension
\citep{trabasso1985causal,graesser1994constructing,kintsch1978macro}.}
The five fields of the situation-model (time, space, protagonist,
causation, intention) are nearly the field set on \texttt{EventNode};
that readers materialise causal links explicitly during
comprehension is the empirical justification for materialising them
explicitly in storage.

\paragraph{Theory of mind
\citep{premack1978tom,wimmer1983false,baroncohen1985autism,
zunshine2006theory,baker2009inverseplanning} and its modern LLM
benchmarks \citep{kim2023fantom,gu2026simpletom,cross2024hyp}.}
\texttt{Belief\{value, confidence, inertia,
established\_at\_fabula, acquired\_via\_event\_id,
acquired\_via\_channel\_id\}} is a discrete approximation of
recursive ToM with provenance. The benchmarks above motivate the
auditor's separate behaviour-given-belief consistency check.

\paragraph{Sternberg's curiosity / suspense / surprise triad
\citep{sternberg1978expositional,sternberg1992telling,
brewer1982stories} and its computational realisations
\citep{cheong2015suspenser,baeyoung2008surprise,wilmot2020suspense,
wilmot2021emnlp,wilmot2022phd}.} The four structural scorers
(\textsc{Mys}, \textsc{Iro}, \textsc{Susp}, \textsc{Sur}) extend the
triad with dramatic irony \citep{booth1974irony,gerrig1993narrative}.
Suspense follows the hope/threat formulation of
\citet{wilmot2020suspense}; surprise follows
\citet{ittibaldi2009surprise} with a geometric prior update; emotional
arc analytics owe to \citet{reagan2016arcs}.

\paragraph{Information-theoretic and epistemic substrate
\citep{shannon1948mathematical,fagin1995reasoning,
alchourron1985agm}.} \texttt{Channel} carries
\texttt{intelligibility} $\in [0,1]$ per participant
(generalising the legacy encryption boolean) and an
\texttt{eavesdropped\_by} relation derived from the
\texttt{intelligibility\_threshold}; belief revision uses an
AGM-style \texttt{beliefs\_added / beliefs\_invalidated} delta.

\paragraph{Causal inference
\citep{pearl1995causal,pearl2009causality,pearl2018book,
spirtes2000causation,bareinboim2022pearl} and actual causality
\citep{halpern2005causes,halpern2016actual,woodward2003making}.}
Pearl's three rungs map directly to the query taxonomy used here
(\texttt{ObservationQuery}, \texttt{InterventionQuery},
\texttt{CounterfactualQuery}); \texttt{mechanism} and
\texttt{causal\_force} on edges are motivated by Halpern's
``actual cause''; ego-graph slicing is a heuristic Markov-blanket
cut \citep{vermapearl1988dsep}. Recent benchmarking of LLMs as
causal reasoners \citep{kiciman2024causal} supports the hybrid
stance taken here: LLMs propose edges (\texttt{extract\_graph.py});
typed code identifies and propagates them
(\texttt{causal\_physics.py}).

\paragraph{Plan-and-render era of neural story generation
\citep{riedlyoung2010narrative,jhalayoung2010cinematic,
martin2018event,yao2019planwrite,goldfarbtarrant2020aristotelian,
rashkin2020plotmachines,yang2022re3,yang2023doc} and its
benchmarks \citep{mostafazadeh2016cloze,fan2018writingprompts,
akoury2020storium,tian2024human,xu2025echoes}.} The
\texttt{CreativeBrief} (Step 9) plus channel-aware draft
(Step 10) is a typed, causal-graph-conditioned variant of the same
plan-then-render pattern. The brief enumerates per-effect typed
\texttt{ConstraintBlock}s; the auditor (Step 11) is the rescore /
revise step \citep{goldfarbtarrant2020aristotelian,yang2022re3}
specialised to narrative-causal rather than stylistic criteria,
in the LLM-as-judge tradition \citep{zheng2023judge,
bai2022constitutional,madaan2023selfrefine,gu2024judge}.

\paragraph{Reading as simulation
\citep{maroatley2008fiction,oatley2016fiction,oatley1995taxonomy,
hogan2003mind}.} The reason a tight causal/affective model
matters more than surface fluency: fiction is a cognitive
simulator, and what is being simulated is a typed world.

\paragraph{Game engines, interactive narrative, and world models
for AI \citep{mateas2003facade,riedl2013interactive,ware2014glaive,
kreminski2020wereyaknow,park2023generative,ha2018world,lecun2022jepa,
hao2023rap,wang2023voyager,smelik2014procedural,
summerville2018pcgml}.} Architecturally, Shadow-Loom is closer to
a game engine than to a sequence model: a typed simulation step
(causal physics) is wrapped by a director / drama-manager
(narrative physics) and exposed to a stochastic policy (the LLM
renderer) only through a constrained action interface, in the
lineage of Mimesis / Fa\c{c}ade-style interactive drama
\citep{mateas2003facade,riedl2013interactive} and narrative
planners such as CPOCL/Glaive \citep{ware2014glaive}. The
generative-agent line of work \citep{park2023generative} and the
``world model'' line in deep RL
\citep{ha2018world,lecun2022jepa,hao2023rap} make the same bet from the opposite direction: that an explicit,
manipulable latent world is what makes long-horizon coherent
behaviour tractable. The shadow branches in Shadow-Loom play the
role those world models play --- a sandbox in which counterfactual
rollouts can be evaluated before being committed --- but with a
typed, inspectable schema in place of an opaque latent state. The bundled example
worlds and authoring tools relate Shadow-Loom to the broader
procedural-content-generation tradition
\citep{smelik2014procedural,summerville2018pcgml,kreminski2020wereyaknow}.


\section{What is novel, what is borrowed}
\label{sec:novelty}

\paragraph{Borrowed.} Every ingredient in
Section~\ref{sec:lineage} is borrowed: Pearl's ladder
\citep{pearl2009causality,bareinboim2022pearl}; the AMWN +
ctf-calculus framework \citep{correa2025amwn}; actual-causality
semantics \citep{halpern2016actual}; the suspense and surprise
formalisms of \citet{wilmot2020suspense} and
\citet{ittibaldi2009surprise}; the LLM-as-judge audit pattern
\citep{bai2022constitutional,madaan2023selfrefine,gu2024judge};
the plan-and-render line in neural story generation
\citep{riedlyoung2010narrative,yao2019planwrite,
goldfarbtarrant2020aristotelian,yang2023doc,yang2022re3,
rashkin2020plotmachines}; the typed-event / story-grammar tradition
\citep{schank1977scripts,lehnert1981plotunits,mani2012computational,
elson2012sig}; and the narratological foundations
(\citealp{shklovsky1965art,genette1980narrative,sternberg1992telling,
ryan1991possible,herman2002story,zwaan1998situation}).

\paragraph{Novel as a combination.} I am not aware of another
open-source pipeline that simultaneously: (a) runs all three
Pearl rungs over a typed narrative graph; (b) tags every node
and edge with an AMWN \texttt{world\_id}
\citep{correa2025amwn} so that counterfactuals are first-class
persisted objects; (c) couples that machinery to closed-form
scorers for the four classical reader-states of narrative
cognition; (d) uses the resulting score as a loss function over
candidate interventions, with the LLM constrained to render the
winner; and (e) closes the loop by re-extracting the rendered
prose and comparing it against the brief. Each ingredient is
borrowed. The composition is the contribution: a typed
two-engine substrate underneath a constrained renderer.
Architecturally the closest analogue is a game engine for
narrative: a typed simulation step gated by a drama-manager-like
scorer
\citep{mateas2003facade,riedl2013interactive,ware2014glaive},
driving a constrained stochastic actor in the spirit of recent
generative-agent and world-model work
\citep{park2023generative,ha2018world,hao2023rap}, but with the
latent state replaced by an explicit causal-graphical schema.

\paragraph{Specific small constructions.} The
``Impact $>$ Inertia'' propagation rule (Eq.~\ref{eq:impact}) and
the geometric prior update inside the surprise scorer
(Eq.~\ref{eq:surprise}) are introduced here as deliberately simple
choices that keep both physics monotonic in useful ways
(App.~\ref{app:narrative}). The two-axis fabula/syuzhet sampling
of the affective scorers (App.~\ref{app:narrative}) is, to my
knowledge, also new.


\section{Rationale for the design choices}
\label{sec:rationale}

\paragraph{Why a typed graph rather than free text + RAG?}
Recent stress tests show that LLMs can answer ``what does X
believe?'' but fail to behave consistently with that belief in
downstream action prediction \citep{gu2026simpletom,kim2023fantom}.
Belief and causal structure that are recoverable from prose are
not, on present evidence, reliably \emph{operated on} inside the
prose-only loop. A typed graph makes belief, channel
intelligibility, and causal mechanism inspectable and
intervenable.

\paragraph{Why both a causal and a narrative physics?}
The causal physics decides \emph{what is allowed to happen} given
the graph; the narrative physics decides \emph{which of the allowed
things would be most readable}. Decoupling licensing from
desirability lets the system score and rank candidate
interventions before any text is generated, and lets the same
engine serve both ``what-if'' analytical queries (rung 2/3) and
high-level affective directives (``maximise dramatic irony for
character X'').

\paragraph{Why constrain the LLM rather than fine-tune one?}
The renderer's job is reduced to dialogue, description, and
pacing within a mathematical envelope. Whatever the model's
training distribution, the brief enumerates the causal edges,
beliefs to preserve, channels to keep hidden, and trait shifts
to honour. Errors that survive into prose are then catchable by
the audit pass because they correspond to specific, testable
claims --- the ``compile and test'' analogue from
Motivation~2.

\paragraph{Why AMWN tagging at the persistence layer?}
Treating a counterfactual as a sibling \texttt{VersionRow} with
\texttt{world\_id="shadow"} makes ``what-if'' branches first-class,
diffable, promotable, and auditable; it is the persistence-layer
analogue of the AMWN graphical construction
\citep{correa2025amwn}. A user can compare a shadow against canon,
keep both, or promote.


\section{Why this might matter beyond fiction}
\label{sec:relevance}

\paragraph{NLP / reasoning.} Narrative is a tractable test bed for
counterfactual and theory-of-mind reasoning because the ground
truth (what happens, who knows what) is finite and authored. A
typed-graph + symbolic-physics + constrained-renderer architecture
is one concrete instantiation of the neuro-symbolic stance for
controllable generation; the audit pass produces a falsifiable
record of where the LLM departed from the brief. Recent ToM
benchmarks \citep{gu2026simpletom,kim2023fantom,cross2024hyp} ask
exactly the kind of behaviour-given-belief questions that an
explicit \texttt{Belief} field with provenance is designed to
support.

\paragraph{Computational social science and digital humanities.}
Many socio-historical and humanistic questions are counterfactual
in form (``what if policy $X$ had not been adopted?'', ``how
would this novel read if the protagonist had spoken sooner?''). A
graphical substrate that lets a researcher ingest a narrative
source, fork a shadow branch under an explicit intervention, and
inspect the propagated trait and relationship deltas is a small
step toward auditable narrative simulation that can sit alongside
established computational-social-science
\citep{lazer2009css} and digital-humanities tooling. I make no
claim of empirical validity for any specific simulation; the
contribution is the substrate, not its calibration.

\paragraph{Author tooling.} The same machinery surfaces in a
NiceGUI workspace and a Model-Context-Protocol server with around
forty tools, exposing the graph, the affective scorers, and the
version DAG to either human authors or LLM agents.

\paragraph{An invitation to other research communities.}
The substrate is deliberately domain-agnostic at the graph level:
entities, events, beliefs, channels, and counterfactual branches
are not specific to fiction. Computational social scientists
\citep{lazer2009css} routinely treat narrative sources --- court
records, parliamentary debates, oral histories, news streams ---
as data, and digital-humanities scholarship has long sought
structured, sharable representations of literary works that can
support both close and distant reading. Shadow-Loom can be read as
one such representation: an open, typed, version-controlled graph
with an explicit counterfactual sandbox and an audit trail back to
the source text. I expressly invite researchers in those
communities, as well as in narratology and cognitive science, to
fork the schema, swap in their own ingestion prompts and affective
scorers, and report what they find. The pipeline ships with
sixteen worked example worlds drawn from canonical literature and
film, which can serve as a common reference set for cross-group
experimentation.


\section{Limitations and future evaluation}
\label{sec:limits}

Shadow-Loom is an experimental research artefact, not a benchmarked
NLP system, and I am explicit about what that does and does not
imply. I report no headline numbers against an external dataset:
the scores produced by the narrative physics are design choices
grounded in narratological theory rather than validated
reader-reaction models, and the causal physics borrows the AMWN
naming and three-rule mental model of \citet{correa2025amwn}
without implementing their identification algorithm literally ---
the graph is heavily typed for narrative use rather than for
rigorous structural causal model identification. The ingestion
pipeline depends on LLM extraction, and although a programmatic
validation and auto-repair loop catches the obvious failure modes,
extraction errors can still propagate. The current unit-test suite
covers the symbolic layers but does not constitute external
evaluation.

Formal evaluation is therefore deferred to future work, along
several complementary axes: (i)~human-judgement studies of the
affective scorers against reader annotations, in the spirit of
\citet{wilmot2020suspense}; (ii)~targeted ablations on bundled
fixtures to quantify the contribution of the causal physics, the
AMWN sandbox, and the audit loop to downstream coherence; and
(iii)~comparison against end-to-end LLM baselines on counterfactual
and theory-of-mind benchmarks \citep{kim2023fantom,gu2026simpletom,
cross2024hyp,xu2025echoes}. I share the present design as a
substrate for further research --- by myself and by others ---
not as a finished product.


\section*{Code, licence, and reproducibility}

The code is released open source under AGPL-3.0-or-later, with a
parallel commercial licence; the authoritative LICENSE file lives
in the project repository. The pipeline, schema, MCP surface,
auditor, and the affective scorers described here all sit in the
\texttt{shadow\_loom/} package. Worked examples on bundled plots
(Macbeth, Death on the Nile, Reservoir Dogs, and others) ship in
\texttt{example\_worlds/} and \texttt{sample\_plots/}.


% \bibliographystyle{acl_natbib}  % set by acl.sty
\bibliography{references}


\appendix
\onecolumn

\section{Definitions and equations}
\label{app:defs}

This appendix gives the formal definitions of the data model and
the equations of each pipeline stage. Section references in the
main paper point here.

\subsection{World model schema (\texttt{WorldStateV1})}
\label{app:schema}

A world state is a tuple
\[
\mathcal{W} = (\mathcal{N}, \mathcal{E}, \mathcal{C}, \mathcal{T}, \tau)
\]
with nodes $\mathcal{N}$ (entities $\mathrm{ENT}\_\bullet$, events
$\mathrm{EVT}\_\bullet$, locations $\mathrm{LOC}\_\bullet$, objects
$\mathrm{OBJ}\_\bullet$, channels $\mathrm{CHN}\_\bullet$, world
traits $\mathrm{WORLD}\_\bullet$), edges $\mathcal{E}$ (causal,
relationship, spatial), channels $\mathcal{C}$ with per-participant
intelligibility $\iota \in [0,1]$, world traits $\mathcal{T}$, and a
branch tag $\tau \in \{\text{factual}, \text{shadow}\}$ inherited
by every node and edge.

Each event $e \in \mathcal{N}_{\mathrm{EVT}}$ carries two integer
indices: chronological $f(e) \in \mathbb{Z}$ (\emph{fabula time})
and presentation $s(e) \in \mathbb{Z}$ (\emph{syuzhet index}), with
$s$ contiguous and unique. Trait values, beliefs, and ambient
state are carried by typed vectors $(v, \iota_{\mathrm{inertia}},
\sigma_{\mathrm{evidence}})$ representing value, inertia, and
evidence strength. Beliefs carry provenance
$(\textit{event\_id}, \textit{channel\_id})$ and the time they
became established. A pure function
$\textsc{Reconstruct}(e_i, t)$ replays sparse
\texttt{state\_timeline} deltas to give the entity's effective
state at any fabula time $t$.

\subsection{Ingestion}
\label{app:ingestion}

Given prose $P$, ingestion produces a $\mathcal{W}$ via:
\begin{enumerate}
\item Five extraction agents over $P$ build a
      \texttt{GlobalRegister} of $\mathrm{WORLD}\_$, $\mathrm{ENT}\_$,
      $\mathrm{LOC}\_$, $\mathrm{OBJ}\_$ nodes plus a fuzzy alias
      table.
\item For each chunk, a Socratic Who/What/Where/When/Why/How
      scaffold dispatches three concurrent specialist agents:
      \emph{physics} (events, causal/spatial edges), \emph{social}
      (relationship edges, channels), \emph{consequences}
      (authoritative entity-state deltas).
\item Output is normalised, fabula times are remapped to a
      uniform $\Delta t = 100$ spacing, and a programmatic
      validator + LLM correction loop enforces ID closure and
      type constraints.
\end{enumerate}

\subsection{Ego-graph slicing and AMWN sandbox}
\label{app:amwn}

Given a query with focal entities $F$ and time anchor $t$, the
ego-graph extractor returns
\[
\mathcal{G}_{F,t} = \big(\textsc{Neighbours}_k(F) \cap
\textsc{Reconstruct}(\cdot, t)\big),
\]
the $k$-hop spatial / causal / social neighbourhood of $F$ at
fabula time $t$, with belief and trait values reconstructed at
$t$ so future information cannot leak into a counterfactual. The
\texttt{AMWNInstantiator} mirrors $\mathcal{G}_{F,t}$ as a
NetworkX \texttt{MultiDiGraph}, in the spirit of the AMWN
construction \citep{correa2025amwn}, with edge types
\{\texttt{located\_in}, \texttt{owned\_by}, \texttt{causal},
\texttt{relationship}, \texttt{connected\_to},
\texttt{communicating\_with}, \texttt{eavesdropped\_by}\}.

\subsection{Causal physics}
\label{app:causal}

Given a sandbox $\mathcal{G}$ and an intervention specification:
\paragraph{Action (rung 2):} $\mathrm{do}(X = x)$ severs every
incoming causal edge into $X$ and writes $V_X \leftarrow x$.
\paragraph{Abduction (rung 3):} for evidence $E$, each entity
trait blends 50\% toward its observed factual value, weighted by
edge $\sigma_{\mathrm{evidence}}$ and gated by a typed
\texttt{MECHANISM\_TRAIT\_MAP} that says which mechanisms can
touch which trait families.
\paragraph{Propagation:} a topological sort over the causal
sub-graph applies, for each downstream node $u$ with parent $p$
and edge weight $w$:
\begin{equation}
\label{eq:impact}
I(u) = (V_p - V_u)\, w,
\qquad
\Delta_u =
\begin{cases}
I(u) - \mathrm{sgn}(I(u))\, \iota_u & \text{if } |I(u)| > \iota_u,\\
0 & \text{otherwise}.
\end{cases}
\end{equation}
A spatial affordance check additionally blocks cross-location
influence unless an unbroken \texttt{connected\_to} path exists.
The result is a structured
\texttt{CausalPhysicsResult\{mutations, blocked, hidden\_deltas,
intervened\_nodes\}}.

\subsection{Narrative physics (affective calculus)}
\label{app:narrative}

Let $A$ be the anchor pair $(t_f, t_s)$ for fabula and syuzhet
cursors, $F$ the focal entities, and write \emph{revealed} for any
event $e$ with $s(e) \le t_s$. Define $\mathrm{Anc}(e)$ as the
causal ancestors of $e$ in the graph and $B(F, t_f)$ as the union
of belief sets of $F$ reconstructed at $t_f$.

\paragraph{Mystery.} For each revealed effect $e$,
\begin{equation}
\label{eq:mystery}
\mathrm{Mys}(e) =
\frac{|\{a \in \mathrm{Anc}(e) : s(a) > t_s\}|}{|\mathrm{Anc}(e)|},
\end{equation}
the fraction of $e$'s causes still hidden from the reader.

\paragraph{Dramatic irony.} For each revealed causal edge
$(a \to e)$ in which $a$ is not in $B(F, t_f)$, count one ``irony
gap''; the score is the fraction of revealed causal connections
that are gaps.

\paragraph{Suspense \citep{wilmot2020suspense}.} For each
unrevealed event $e$ with strongest incoming causal weight $p_e$:
\begin{equation}
w_{\mathrm{threat}} = \!\!\sum_{e:\, F \subseteq \mathrm{target}(e)} \!\!p_e,
\quad
w_{\mathrm{hope}} = \!\!\sum_{e:\, F \cap \mathrm{actor}(e) \neq \emptyset} \!\!p_e,
\end{equation}
\begin{equation}
\mathrm{Susp} = \mathrm{clip}_{[0,1]}\!\left(
\frac{w_{\mathrm{threat}} - w_{\mathrm{hope}}}
     {w_{\mathrm{threat}} + w_{\mathrm{hope}}}
\right),
\end{equation}
returning $0$ when $w_{\mathrm{hope}} = 0$ (the
suspense--despair boundary).

\paragraph{Surprise.} Per trait $T$, the prior $q$ starts at the
corpus marginal across entities (defaulting to $0.5$), then for
each revealed cause $(a \to e)$ with weight $w$ it is updated by
\begin{equation}
\label{eq:surprise}
q \mathrel{+}= w\,(\mathrm{actual}_T - q),
\end{equation}
a geometric pull that monotonically converges on the truth.
Letting $p$ be the actual trait value,
\begin{equation}
D_{\mathrm{KL}}(p \,\|\, q) =
p \log \tfrac{p}{q} + (1-p)\log \tfrac{1-p}{1-q},
\end{equation}
and the entity-level surprise is the mean over traits, normalised
by $\log(1/\varepsilon)$ to $[0,1]$ for a small $\varepsilon$.

\paragraph{Six emotion scorers} (\texttt{grief, rage, joy, regret,
love, fear}) declare which traits should move in which direction
and score the headroom available, weighted by inertia evidence.
The combined affective score requested by a
\texttt{DirectiveQuery} is a weighted sum over the active
functionals.

\paragraph{Two-axis sampling.} The same scorers are evaluated
along the syuzhet axis (advancing $t_s$ only reveals more, so
mystery, irony, surprise are non-increasing) and the fabula axis
(snapshotting the world at $t_f$, where surprise can spike at
intrinsically surprising story moments).

\subsection{Generation}
\label{app:generation}

\texttt{DirectiveAssembler.evaluate\_candidate\_events()} forks the
sandbox per candidate intervention, runs the causal physics, prunes
candidates that violate inertia, affordance, or propagation
constraints, ranks survivors by the affective score, and packages
the winner as a \texttt{CreativeBrief} of typed
\texttt{ConstraintBlock}s (per-effect): hidden-predecessor lists for
mystery, ``MUST NOT learn'' guards for irony, threat/hope tables and
hidden-channel lists for suspense, per-trait KL magnitudes for
surprise, and per-trait shift constraints with headroom and inertia
evidence for the emotion targets. Only the brief is shown to the
renderer LLM, which produces dialogue, description, and pacing
within the envelope.

\subsection{Audit and feedback loop}
\label{app:audit}

The auditor \citep[an LLM-as-judge in the sense of][]{gu2024judge}
runs three passes over the rendered prose:
(i)~\textbf{causal} --- reverse-engineers prose into causal claims
and flags ``miracle steps'' (state changes with no licensing edge
in the brief); (ii)~\textbf{abduction} --- runs counterfactual
probes against the prose to check that implicit events hold;
(iii)~\textbf{affective} --- measures the realised epistemic gap
in the prose against the target. If the structured loss is
non-zero, generation is re-run with the auditor's feedback
appended, up to \texttt{max\_correction\_retries}. On success the
re-extracted prose is merged into a new \texttt{VersionRow}, with
\texttt{world\_id} set by the branch policy
($\texttt{auto} \mid \texttt{mainline} \mid \texttt{shadow}$);
counterfactual queries default to a fresh \texttt{shadow} branch,
preserving canon.


\section{Worked examples on real plot fixtures}
\label{app:examples}

The repository ships sixteen hand-authored \texttt{example\_worlds/}
fixtures paired with prose synopses in \texttt{sample\_plots/}. A
small subset is used below to illustrate each component
end-to-end. Field
values quoted below are verbatim from the fixtures; queries are the
same ones run by the example walkthroughs in
\texttt{docs/model-examples.md} and \texttt{docs/pipeline-by-example.md}.

\subsection{World model --- Macbeth}
\label{app:ex-worldmodel}

In \texttt{example\_worlds/macbeth.py}, the location
\texttt{LOC\_HEATH} carries
\texttt{ambient\_state\{supernatural: AmbientVector(value=0.9,
volatility=0.2, evidence\_strength="strong"), concealment: 0.7\}}.
The entity \texttt{ENT\_MACBETH} is a bundle of typed
\texttt{TraitVector}s: \texttt{ambition}, \texttt{courage} (high
inertia $0.7$, hard to shake), \texttt{guilt} (low inertia $0.25$,
highly mutable), and so on; an initial
\texttt{Belief(target=PROPHECY, confidence=0.4)} sits on his belief
list with its own inertia. \texttt{OBJ\_BLOODY\_DAGGERS} declares
\texttt{Affordance(action="kill", target\_type="Entity")} ---
the precondition gate the engine checks before it lets
\texttt{EVT\_DUNCAN\_MURDER} fire. Every \texttt{EventNode} carries
both indices: the murder has a fabula time but the syuzhet index
follows the staged scene. The fixture exercises all five
\texttt{CausalEdge} modalities (\texttt{chain\_reaction},
\texttt{mutation}, \texttt{mutation\_social},
\texttt{affordance\_gate}, \texttt{ambient\_propagation}); the
\texttt{ambient\_propagation} edge is what lets
\texttt{LOC\_HEATH.supernatural=0.9} bias Macbeth's uptake of the
prophecy belief.

\subsection{AMWN sandbox --- Macbeth ``what if Macbeth refuses?''}
\label{app:ex-amwn}

The query ``Macbeth refuses to murder Duncan'' is a rung-2
intervention. \texttt{AMWNInstantiator.create\_sandbox(F=\{ENT\_MACBETH\},
t)} mirrors the ego-graph as a \texttt{MultiDiGraph}; a sibling
\texttt{VersionRow} with \texttt{world\_id="shadow"} holds the fork.
\texttt{apply\_do\_operator(\{EVT\_DUNCAN\_MURDER: $\oslash$\})}
severs the murder's incoming causal edges; the downstream
\texttt{chain\_reaction} edges into \texttt{EVT\_MALCOLM\_FLEES},
\texttt{EVT\_MACBETH\_CROWNED}, \texttt{EVT\_BANQUO\_MURDERED} all
lose activation; the \texttt{mutation} edges into
\texttt{ENT\_MACBETH.guilt} and \texttt{paranoia} never fire. The
factual mainline is untouched --- the shadow row is browsable,
diffable, and promotable from the UI.

\subsection{Causal physics --- Macbeth propagation under
``Impact \texorpdfstring{$>$}{>} Inertia''}
\label{app:ex-causal}

After the do-operator, \texttt{propagate()} topologically sorts the
causal sub-graph. For the original (factual) chain, the edge
\texttt{EVT\_BANQUO\_GHOST $\to$ ENT\_MACBETH.guilt} with weight
$w=0.7$ and source delta $V_p - V_u = 0.6$ produces
$I = 0.42$. Because \texttt{guilt}'s post-shock inertia $\iota$ has
been bumped to $\sim 0.3$ from its baseline $0.25$ by the previous
\texttt{state\_timeline} entry, the gate $|I| > \iota$ passes and
$\Delta = 0.42 - 0.3 = 0.12$ is written. A symmetrical attempt to
shock \texttt{courage} (initial inertia $0.7$) at $w=0.5$,
$\Delta_\text{src}=0.4$ gives $I=0.20<0.7$ and is \emph{blocked} ---
exactly the play's psychology, in which his nerve outlasts his
conscience. Spatial affordance gating additionally blocks
cross-location shocks unless a \texttt{connected\_to} path exists.

\subsection{Narrative physics --- Death on the Nile, Reservoir
Dogs, Romeo and Juliet, Macbeth}
\label{app:ex-narrative}

The four structural scorers each have a canonical fixture.

\paragraph{Mystery (Macbeth Act II).} Eq.~\ref{eq:mystery} on the
revealed effect \texttt{EVT\_DUNCAN\_FOUND\_DEAD}: of its causal
ancestors, \texttt{EVT\_LADY\_PERSUADES} and
\texttt{EVT\_DAGGERS\_PLACED} are still hidden at the audience's
\texttt{syuzhet\_anchor} early in the act; the score is
$2/3 \approx 0.67$, falling to $0$ once Act V resolves the
provenance.

\paragraph{Dramatic irony (Death on the Nile).} The reader knows
Jacqueline and Simon are conspirators (\texttt{EVT\_PLAN\_HATCHED}
revealed at low syuzhet); Poirot does not. For each later event
$(a \to e)$ where $a$ is the conspiracy and $a \notin
B(\text{ENT\_POIROT}, t_f)$, an irony gap is counted. The score
sits high through the cruise scenes and only collapses at the
denouement.

\paragraph{Suspense (Romeo and Juliet, tomb scene).} At the tomb
$F=\{\text{ENT\_ROMEO}\}$, unrevealed events with Romeo as
non-acting target (Friar's letter intercepted, Juliet's draught
about to wear off) sum to $w_\text{threat}$; events Romeo himself
authors and that have not yet fired (drink poison, kill Paris) sum
to $w_\text{hope}$. Threat dominates. As Romeo drinks, $w_\text{hope}
\to 0$ and the score returns $0$ --- the Wilmot--Keller
suspense--despair boundary is hit cleanly.

\paragraph{Surprise (Reservoir Dogs reveal).}
\texttt{EVT\_ORANGE\_RECRUITED} sits at low fabula time but high
syuzhet; before the reveal the corpus-marginal prior on
\texttt{ENT\_ORANGE.role} starts at $q=0.5$, and the few revealed
edges into Orange give a small geometric pull keeping $q \approx
0.6$ ``criminal''. After the reveal the actual value $p \approx
0.95$ ``cop''. Eq.~\ref{eq:surprise} gives a per-trait KL spike
that is large but \emph{licensed} --- the engine had the cause
present at low fabula time, only locked behind syuzhet, so
re-extraction will not flag a miracle step.

\paragraph{Emotion targets (Gone Girl, Brief Encounter).} The six
trait-trajectory scorers (\texttt{grief, rage, joy, regret, love,
fear}) declare which traits should move in which direction.
On Gone Girl, the directive ``maximise grief in Nick'' resolves to
required downshifts on \texttt{trust} and \texttt{hope} on
\texttt{ENT\_NICK}, weighted by inertia evidence; the available
headroom score determines whether the directive is even feasible
at the chosen anchor. On Brief Encounter the same machinery
realises \texttt{regret} as Laura's
$\Delta(\texttt{self\_actualisation})$ headroom under a
constant-supervision channel that never permits the choice to be
made.

\subsection{Directive assembly --- Romeo and Juliet}
\label{app:ex-directive}

The directive ``raise dramatic irony to $0.85$ in the tomb scene
without giving Romeo correct information'' triggers
\texttt{DirectiveAssembler.evaluate\_candidate\_events()}. From
unblocked affordances and reachable spatial paths it enumerates
candidates --- \emph{Friar's letter intercepted}, \emph{Balthasar
arrives faster}, \emph{Juliet wakes one minute earlier} --- and
forks an AMWN sandbox per candidate. The plausibility gate drops
\emph{Balthasar arrives faster} (the plague quarantine
\texttt{SpatialEdge.is\_locked} blocks the path); the remaining
sandboxes are propagated, and the survivors are scored. \emph{Friar's
letter intercepted} wins because it raises Eq.~\ref{eq:mystery}'s
sister gap-count for irony without adding any utterance event into
$B(\text{ENT\_ROMEO}, t_f)$. The winner is wrapped in typed
\texttt{ConstraintBlock}s: a \texttt{must\_events} list, a
\texttt{must\_not\_events} list (Romeo learning the truth), trait
envelopes (Romeo's \texttt{despair} envelope), and a syuzhet-window
constraint. That bundle is the \emph{only} thing handed to the
renderer.

\subsection{Generation --- the brief as safety envelope}
\label{app:ex-generation}

The renderer LLM is given the Romeo-and-Juliet brief as system
prompt plus a small style context. It may write any prose about
the tomb, the messenger, the dagger; it may not invent a causal
edge, shift a trait outside its envelope, open a channel that does
not exist, or place an utterance the brief does not licence. The
output is a scene of dialogue, description, and pacing inside the
mathematical envelope chosen by directive assembly --- a typed,
causal-graph-conditioned instance of the plan-and-render pattern
\citep{riedlyoung2010narrative,yao2019planwrite,
goldfarbtarrant2020aristotelian,yang2023doc}.

\subsection{Audit and feedback loop --- Gone Girl}
\label{app:ex-audit}

Gone Girl is the system's lying-narrator fixture: Amy's diary is a
sequence of \texttt{EventNode}s with \texttt{truth\_value="false"}
broadcast through \texttt{CHN\_DIARY\_PUBLIC}. After generation,
\texttt{auditor.run\_feedback\_loop} runs three passes in parallel.
The \emph{causal audit} reverse-engineers the prose into
$(\text{source}, \text{target}, \text{modality})$ claims and
matches each against the brief; an LLM-emitted ``Nick suddenly
realises Amy is alive'' with no licensing causal edge is flagged
as a \emph{miracle step}. The \emph{abduction audit} runs
counterfactual probes (``if the diary were truthful, would the
prose still hold?''); a ``yes'' here means the prose has
collapsed Amy's two faces into one, and is rejected. The
\emph{affective audit} measures the realised dramatic-irony gap
against the target. The aggregated \texttt{FeedbackLoopResult}
exposes both an \texttt{engine\_thresholds\_passed} gate
(deterministic: number of miracle steps, max trait drift) and an
\texttt{llm\_pass} gate (literary critique); either failing
re-runs generation with the audit feedback appended, up to
\texttt{max\_correction\_retries}. On convergence the scene is
re-extracted into a new \texttt{VersionRow}; counterfactual
branches land on \texttt{world\_id="shadow"} and only become canon
via explicit \texttt{promote\_branch} \citep{correa2025amwn}.

\end{document}
